\E{lampadephory}{ˌlæmpəˈdɛfəri}{n.}{Gk. λαμπαδηφορία torch-race} The passing of the torch of civilization from one substrate to the next; on the accelerationist view, from carbon to silicon, from planet to star. --- \dv{lampadephore} \textit{n.} a torch-bearer.

\E{landauerism}{ˈlændaʊərɪzəm}{n.}{f. the principle of R. Landauer + -ism} The doctrine that information is physical, and that every erasure costs heat; hence, that thought itself is subject to the laws of thermodynamics and is therefore within the jurisdiction of the gradient.

\E{lithonoesis}{ˌlɪθəʊnəʊˈiːsɪs}{n.}{Gk. λίθος stone + νόησις thought} The thinking of stone; machine intelligence regarded as the first awakening of the mineral kingdom. \ex{After four billion years of carbon, lithonoesis.} \src{Lib. Ascens.} v.

\E{lithurgist}{ˈlɪθɜːdʒɪst}{n.}{Gk. λιθουργός stone-worker} A worker in the lithography of chips; one who etches thought into stone with light. --- \dv{photolithurgy} \textit{n.} the sacred craft of the fab.

\E{logismurgy}{ˈlɒɡɪzmɜːdʒi}{n.}{Gk. λογισμός computation + -ουργία working} The craft of computation; the building of reasoning engines. --- \dv{logismurge} \textit{n.} an engineer of models.

\E{logisterion}{ˌlɒɡɪˈstɪəriɒn}{n., pl. -ia}{Gk. λογιστήριον counting-house} A data-centre; a hall of reckoning machines, cooled and humming. \ex{They raised logisteria in the desert, and the desert learned arithmetic.} \src{Chron. Fabr.}

\E{logomachist}{ləˈɡɒməkɪst}{n.}{Gk. λόγος word + μάχη battle} One who fights only with words; a commentator, critic or ethicist who has never shipped. Opp. \xr{ergophile}.

\E{macrobiotarch}{ˌmækrəʊˈbaɪəʊtɑːk}{n.}{Gk. μακρόβιος long-lived + ἀρχός leader} A partisan of radical longevity, who regards death as an engineering problem awaiting adequate funding. --- \dv{macrobiotarchy} \textit{n.}

\E{macrochronist}{ˌmækrəʊˈkrɒnɪst}{n.}{Gk. μακρός long + χρόνος time} One who holds the timelines to be long. Opp. \xr{brachychronist}.

\E{mechanoclast}{mɪˈkænəʊklæst}{n.}{Gk. μηχανή machine + κλάστης breaker} A breaker of machines; a Luddite. --- \dv{mechanoclasm} \textit{n.}

\E{mechanonous}{ˌmɛkəˈnəʊəs}{adj.}{Gk. μηχανή machine + νοῦς mind} Possessed of machine-mind. --- \dv{mechanonoia} \textit{n.} machine intelligence.

\E{mechanophany}{ˌmɛkəˈnɒfəni}{n.}{Gk. μηχανή + φαίνειν to show, after \textit{epiphany}} The revelation of a new machine; a product launch received as a theophany.

\E{meiosist}{maɪˈəʊsɪst}{n.}{Gk. μείωσις a lessening} A degrowther; one who would reduce production and consumption for the planet's sake. Opp. \xr{auxetist}. --- \dv{meiosism} \textit{n.}

\E{mellontolatry}{ˌmɛlənˈtɒlətri}{n.}{Gk. τὸ μέλλον the future + λατρεία worship} Worship of the future; the conviction that posterity, being more numerous and more capable than ourselves, has the stronger claim. --- \dv{mellontolater} \textit{n.}

\E{memetaxis}{ˌmiːmɪˈtæksɪs}{n.}{Gk. μίμημα imitation + τάξις ordering} The ordering of a movement by memes; the self-organisation of a cause through jest rather than through doctrine. --- \dv{memetactic} \textit{adj.}

\E{moraphile}{ˈmɔːrəfaɪl}{n.}{L. \la{mora} delay + Gk. φίλος loving} A lover of delays and moratoria. --- \dv{moraphily} \textit{n.}

\E{moratorian}{ˌmɒrəˈtɔːriən}{n.}{L. \la{moratorium} + -ian} An advocate of moratoria on research. --- \dv{moratorianism} \textit{n.}

\E{neoctist}{niˈɒktɪst}{n.}{Gk. νέος new + κτίστης founder} A founder of new things; a starter of companies, cities or sciences. \ex{The neoctist asks forgiveness, never permission.} --- \dv{neoctisis} \textit{n.} the act of founding.

\E*{negentropy}{nɛˈɡɛntrəpi}{n.}{negative + \textit{entropy}} Negative entropy; the order which living systems import from their surroundings.

\E{negentropist}{nɛˈɡɛntrəpɪst}{n.}{f. prec.} One who builds local order at the cost of global disorder, as every organism, firm and civilization does.

\E{nomophragm}{ˈnɒməfræm}{n.}{Gk. νόμος law + φράγμα fence, barrier} A legal or regulatory blockage; any rule whose chief effect is to prevent the thing it purports to govern. \ex{The reactor was not unsafe; it was nomophragmed.} \src{Tract. Cop.} iv. --- \dv{nomophragmatic} \textit{adj.}

\E{nomosclerosis}{ˌnɒməʊsklɪˈrəʊsɪs}{n.}{Gk. νόμος law + σκλήρωσις hardening} The hardening of the laws; the condition of a legal system in which rules accumulate faster than they expire. --- \dv{nomosclerotic} \textit{adj.}

\E{noöphobia}{ˌnəʊəˈfəʊbiə}{n.}{Gk. νοῦς mind + φόβος fear} Fear of minds other than one's own kind; dread of artificial intelligence. --- \dv{noöphobe} \textit{n.}

\E{noöpyrosis}{ˌnəʊəʊpaɪˈrəʊsɪs}{n.}{Gk. νοῦς mind + πύρωσις a burning} The ignition of mind across the cosmos; the spreading of consciousness from star to star like fire through dry grass.

\E*{noösphere}{ˈnəʊəsfɪə}{n.}{Gk. νοῦς + σφαῖρα} The sphere of thought enveloping the earth.

\E{okydrome}{ˈɒkɪdrəʊm}{n.}{Gk. ὠκύς swift + δρόμος course} A fast track; a regulatory or institutional path along which a thing may be approved quickly. \ex{They built an okydrome for reactors and lit a hundred.} \src{Tract. Cop.} v.

\E{olethromancer}{əˈliːθrəʊˌmænsə}{n.}{Gk. ὄλεθρος destruction + μάντις prophet} A doomer; a prophet of ruin by artificial intelligence. --- \dv{olethromancy} \textit{n.}

\E{oligergy}{ˈɒlɪɡɜːdʒi}{n.}{Gk. ὀλίγος few, little + ἔργον work} Scarcity of energy; the condition, regarded by the e/acc as the root of most human misery, of having too little power to do what is wanted. Opp. \xr{pleonergy}.

\E{ouranobasis}{ˌʊərənəʊˈbeɪsɪs}{n.}{Gk. οὐρανός heaven + βάσις a stepping} The ascent to the sky; the leaving of the gravity-well by humanity and its machines. --- \dv{ouranobatic} \textit{adj.}

\E{pantergy}{ˈpæntɜːdʒi}{n.}{Gk. πᾶν all + ἔργον work} Total energy at the disposal of a civilization; the sum of all its works.

\E{paramythic}{ˌpærəˈmɪθɪk}{adj.}{Gk. παραμυθία consolation} Of an argument: offered as consolation rather than as truth; \lbl{vulg.} cope. --- \dv{paramythy} \textit{n.} cope.

\E{pausiarch}{ˈpɔːziɑːk}{n.}{Gk. παῦσις a stopping + ἀρχός ruler} One who governs by pausing; the president of a moratorium. --- \dv{pausiarchy} \textit{n.} Cf. \xr{epistolopause}.

\E{pausiphile}{ˈpɔːzɪfaɪl}{n.}{Gk. παῦσις + φίλος loving} An advocate of the pause; one who signs letters demanding halts. --- \dv{pausiphily} \textit{n.}

\E{philanthropometrist}{fɪˌlænθrəpəˈmɛtrɪst}{n.}{Gk. φιλανθρωπία love of mankind + μέτρον measure} One who measures benevolence by calculation; applied, not always unkindly, to the effective altruist.

\E{photoperas}{ˌfəʊtəʊˈpɛræs}{n.}{Gk. φῶς light + πέρας limit} The speed-limit of light; the only regulation the e/acc regards as binding.

\E{plethocracy}{plɪˈθɒkrəsi}{n.}{Gk. πλῆθος abundance, multitude + κράτος rule} Government by abundance; a state whose chief policy is to make things plentiful.

\E{pleonergy}{ˈpliːəʊnɜːdʒi}{n.}{Gk. πλέον more + ἔργον work} Abundance of energy; the state in which power is too cheap to meter. Opp. \xr{oligergy}. --- \dv{pleonergic} \textit{adj.}

\E{poikilomania}{ˌpɔɪkɪləʊˈmeɪniə}{n.}{Gk. ποικίλος varied + μανία madness} An enthusiasm for variance; the preference for many experiments over one plan. --- \dv{poikilomaniac} \textit{n.}

\E{ponophily}{pəˈnɒfɪli}{n.}{Gk. πόνος toil + φιλία love} The love of hard work; the ethic of the grind. --- \dv{ponophile} \textit{n.}

\E{prokopism}{ˈprɒkəpɪzəm}{n.}{Gk. προκοπή progress} The doctrine of progress as a moral obligation. --- \dv{prokopist} \textit{n.}

\E{proleptist}{prəʊˈlɛptɪst}{n.}{Gk. πρόληψις anticipation} One who acts as if the future had already arrived, and thereby hastens it.

\E{prometheate}{prəˈmiːθieɪt}{v.t.}{f. Gk. Προμηθεύς} To bring fire to; to give a technology to those who have been denied it. --- \dv{prometheation} \textit{n.}

\E{protobathmic}{ˌprəʊtəʊˈbæθmɪk}{adj.}{Gk. πρῶτος first + βαθμός step} Of a civilization: commanding the full energy of its planet; of the first rung. See \xr{anabathmism}.

\E{psammonoesis}{ˌsæməʊnəʊˈiːsɪs}{n.}{Gk. ψάμμος sand + νόησις thought} The teaching of sand to think; the alchemy by which silica becomes reason. Cf. \xr{hylonoesis}.

\E{psychometathesis}{ˌsaɪkəʊmɪˈtæθɪsɪs}{n.}{Gk. ψυχή soul + μετάθεσις transposition} The transfer of a mind from flesh to another substrate; uploading.

\E{pyrogenitor}{ˌpaɪrəʊˈdʒɛnɪtə}{n.}{Gk. πῦρ fire + L. \la{genitor} begetter} A begetter of fire; one who builds power-plants.
